Deux ampoules sont conçues pour fonctionner sous \qty{120}{V} : l'une a une puissance nominale de \qty{\v{P1}}{W}, l'autre de \qty{\v{P2}}{W}. On considère que leur résistance est constante.

\begin{enumerate}[label=\alph*)]
    \item Déterminez la résistance de chaque ampoule.
    \item On branche les deux ampoules en série sur une prise de \qty{120}{V}. Déterminez l'intensité du courant dans le circuit.
    \item Déterminez la puissance réellement dissipée par chaque ampoule dans ce montage en série. Laquelle brille le plus ?
\end{enumerate}
